% !TeX root = ../../main.tex

Die Diagramme~\ref{fig:results} beschreiben die wichtigsten Kenngrößen über die durchgeführten Experimente und
ermöglichen einen Vergleich der drei untersuchten Fehlerklassen.
Bei der Erfolgsrate zeigen sich merkliche Unterschiede zwischen den Fehlerklassen.
Deadlocks erreichen über alle Experimente eine durchschnittliche Erfolgsrate von 45.66\%.
Bei Lock Timeout Fehlern liegt diese bei 13.69\%, während Serialization Fehler mit 56.33\% die einzige Fehlerklasse
mit einer Erfolgsrate von über 50\% darstellen.

Auch bei der Ausführungszeit zeigen sich deutliche Unterschiede.
Erfolgreiche Transaktionen weisen bei Deadlocks mit durchschnittlich 4602.26ms die höchste Ausführungszeit auf.
Bei Lock Timeout Fehlern beträgt diese 1232.09ms und bei Serialization Fehlern lediglich 165.11ms.
Ein ähnliches Verhältnis zeigt sich bei den fehlgeschlagenen Transaktionen.
Hier benötigen Deadlocks mit durchschnittlich 10739.17ms ebenfalls länger, bevor die Transaktion
abgebrochen wird.
Bei Lock Timeout Fehlern beträgt die durchschnittliche Ausführungszeit 509.93ms und bei Serialization Fehlern 214.39ms.

Analog dazu ist der Retry Overhead.
Hier entsteht bei Deadlock Fehlern passend zu der Ausführungszeit sowohl bei erfolgreichen als auch bei
fehlgeschlagenen Transaktionen der größte zusätzliche Zeitaufwand.
Bei einer erfolgreichen Transaktion liegt der Overhead bei 3511.47ms und bei einer fehlgeschlagenen bei 8854.39ms.
Auch hier ist der Overhead bei Lock Timeouts höher als der bei Serialization Fehlern, wie es bei der
Ausführungszeit schon der Fall ist.
Ein Lock Timeout hat bei einer erfolgreichen Transaktion 957.19ms und bei einer fehlgeschlagener Transaktion
408.65ms an Overhead erzeugt.
Serialization Fehler verursachen dagegen mit 123.16ms bei erfolgreichen und 170.44ms bei fehlgeschlagenen
Transaktionen einen geringeren Overhead.

Beim Durchsatz zeigen sich ebenfalls erkennbare Unterschiede zwischen den Fehlerklassen.
Über alle Experimente weisen Deadlocks im Durchschnitt mit 0.5185 Transaktionen pro Sekunde den
niedrigsten Durchsatz auf.
Innerhalb dieser Fehlerklasse erreicht \texttt{Experiment1} mit 0.5604tx/s den höchsten Wert, während
\texttt{Experiment2} mit 0.4627tx/s den niedrigsten Durchsatz aufweist.
Der durchschnittliche Durchsatz bei Lock Timeouts liegt mit 0.9923 Transaktionen pro Sekunde höher.
Der höchste Wert wird in \texttt{Experiment2} mit 0.9934tx/s erreicht während \texttt{Experiment0} mit 0.9903tx/s
den niedrigsten Durchsatz aufweist.
Die Unterschiede zwischen den einzelnen Experimenten fallen hier somit geringer aus als bei den Deadlocks.
Serialization-Fehler weisen mit durchschnittlich 15.0225tx/s den mit Abstand höchsten Durchsatz der drei
Fehlerklassen auf.
Der höchste Wert wird in \texttt{Experiment1} mit 16.3201tx/s erreicht.
Der niedrigste Durchsatz tritt mit 13.9250tx/s in \texttt{Experiment0} auf.

% !TeX root = ../main.tex

\begin{figure}[ht]
	\centering
	\begin{minipage}{0.45\textwidth}
		\centering
		\caption*{(a) Erfolgs- und Fehlerraten}
		% !TeX root = ../../main.tex

\begin{tikzpicture}
	\begin{axis}[
		xbar,
		bar width=2pt,
		enlarge y limits=0.3,
		width=0.7\textwidth,
		height=0.5\textwidth,
		trim axis left,
		trim axis right,
		xmin=0,
		xmax=100,
		xlabel={\%},
		symbolic y coords={Deadlock, Serialization, Lock Timeout},
		ytick=data,
		legend style={
			at={(0.5, -0.75)},
			anchor=north,
			legend columns=2
		},
	]

		% Success
		\addplot[fill=successcolor, draw=none] coordinates {
			(45.66,Deadlock)
			(56.33,Serialization)
			(13.69,Lock Timeout)
		};

		% Failed
		\addplot[fill=failedcolor, draw=none] coordinates {
			(54.34,Deadlock)
			(43.67,Serialization)
			(86.31,Lock Timeout)
		};

		\legend{Success, Failed}

	\end{axis}
\end{tikzpicture}
	\end{minipage}
	\begin{minipage}{0.45\textwidth}
		\centering
		\caption*{(b) Ausführungszeiten}
		% !TeX root = ../../main.tex

\begin{tikzpicture}
	\begin{axis}[
		xbar,
		bar width=2pt,
		enlarge y limits=0.3,
		width=0.7\textwidth,
		height=0.5\textwidth,
		xmin=0,
		xmax=11000,
		xlabel={ms},
		scaled ticks=false,
		symbolic y coords={Deadlock, Serialization, Lock Timeout},
		ytick=data,
		legend style={
			at={(0.5,-0.75)},
			anchor=north,
			legend columns=2
		},
	]

		% Success
		\addplot[fill=successcolor, draw=none] coordinates {
			(4602.26,Deadlock)
			(165.11,Serialization)
			(1232.09,Lock Timeout)
		};

		% Failed
		\addplot[fill=failedcolor, draw=none] coordinates {
			(10739.17,Deadlock)
			(214.39,Serialization)
			(509.93,Lock Timeout)
		};

		\legend{Success, Failed}

	\end{axis}
\end{tikzpicture}
	\end{minipage}
	\begin{minipage}{0.45\textwidth}
		\centering
		\caption*{(c) Overhead}
		% !TeX root = ../../main.tex

\begin{tikzpicture}
	\begin{axis}[
		xbar,
		bar width=2pt,
		enlarge y limits=0.3,
		width=0.7\textwidth,
		height=0.5\textwidth,
		trim axis left,
        trim axis right,
		xmin=0,
		xmax=11000,
		xlabel={ms},
		scaled ticks=false,
		symbolic y coords={Deadlock, Serialization, Lock Timeout},
		ytick=data,
		legend style={
			at={(0.5,-0.75)},
			anchor=north,
			legend columns=2
		},
	]

		% Success
		\addplot[fill=successcolor, draw=none] coordinates {
			(3511.47,Deadlock)
			(123.16,Serialization)
			(957.19,Lock Timeout)
		};

		% Failed
		\addplot[fill=failedcolor, draw=none] coordinates {
			(8854.39,Deadlock)
			(170.44,Serialization)
			(408.65,Lock Timeout)
		};

		\legend{Success, Failed}

	\end{axis}
\end{tikzpicture}
	\end{minipage}
	\begin{minipage}{0.45\textwidth}
		\centering
		\caption*{(d) Durchsatz}
		% !TeX root = ../../main.tex

\begin{tikzpicture}
	\begin{axis}[
		xbar,
		bar width=2pt,
		enlarge y limits=0.3,
		width=0.7\textwidth,
		height=0.5\textwidth,
		xmin=0,
		xmax=17,
		xlabel={tx/s},
		scaled ticks=false,
		symbolic y coords={Deadlock, Serialization, Lock Timeout},
		ytick=data,
		legend style={
			at={(0.5,-0.75)},
			anchor=north,
			legend columns=4
		},
	]

		% e0
		\addplot[fill=ezerocolor, draw=none] coordinates {
			(0.5485,Deadlock)
			(13.9250,Serialization)
			(0.9903,Lock Timeout)
		};

		% e1
		\addplot[fill=eonecolor, draw=none] coordinates {
			(0.5604,Deadlock)
			(16.3201,Serialization)
			(0.9926,Lock Timeout)
		};

		% e2
		\addplot[fill=etwocolor, draw=none] coordinates {
			(0.4627,Deadlock)
			(14.7587,Serialization)
			(0.9934,Lock Timeout)
		};

		% e3
		\addplot[fill=ethreecolor, draw=none] coordinates {
			(0.5024,Deadlock)
			(15.0863,Serialization)
			(0.9930,Lock Timeout)
		};

		\legend{\(\mathrm{e}_0\), \(\mathrm{e}_1\), \(\mathrm{e}_2\), \(\mathrm{e}_3\)}

	\end{axis}
\end{tikzpicture}
	\end{minipage}

	\caption{Wichtige Messgrößen}
	\label{fig:results}
\end{figure}